\section{Related Work}\label{sec:related}
\subsection{Classical link adaptation and OLLA}
Outer-loop link adaptation (OLLA) corrects the reported channel quality indicator (CQI) with an offset that steps down after a NACK and up after an ACK, so that the long-run block-error rate approaches a target~\cite{nakamura2002olla,pedersen2007fdps}. Variants adapt the target~\cite{park2015targetbler}, the step sizes~\cite{blanquez2016eolla,sarret2015dolla,ohseki2016fastolla}, the offset estimator~\cite{delgado2017seqtest,zhu2023nolla}, or the feedback used under URLLC soft-ACK~\cite{peralta2022ollaurllc}. Ahsen et al.\ use a rank-dependent multi-offset rule~\cite{ahsen2026moolla}; this is a different design from the Multi-OLLA comparator used here. SALAD varies three OLLA step sizes without learning~\cite{wiesmayr2025salad}.

\subsection{Learned link adaptation}
Bandit formulations choose the modulation and coding scheme (MCS) or the offset directly~\cite{saxena2019cmab,saxena2022tuningfree,park2021twostagets,ku2024cmaburllc}. Deep reinforcement learning (RL) has been applied with DQN, actor-critic, TD3 and PPO agents~\cite{ye2023drlla,pang2024td3olla,you2026sim2real,parsa2025qdrlla,demirel2025genrl}. Deadline-aware RL on an OLLA loop is studied by Kela et al.~\cite{kela2022delayolla}, and learning on top of OLLA by Pulliyakode and Kalyanasundaram~\cite{pulliyakode2017rlolla}, Pang et al.~\cite{pang2024td3olla} and Wang et al. (LOLLA)~\cite{wang2026lolla}. These works are evaluated in different settings from ours (LOLLA, for example, treats uplink with an immediate CRC outcome and no HARQ feedback delay), and we do not re-implement them. In O-RAN, learning-based control is placed by timescale: xSlice, a PPO agent with a graph convolutional network, runs as an xApp on the near-real-time RIC, whose control loop spans 10~ms to 1~s, and re-partitions resource blocks between slices over the air~\cite{yan2026xslice}, whereas LOLLA runs as a dApp~\cite{doro2022dapps} inside the GPU-accelerated stack and sets the MCS in every slot, with end-to-end control latency below 500~$\mu$s at 30~kHz subcarrier spacing~\cite{wang2026lolla}. Our question is a different one: given a simulator with multi-process HARQ, how much of a learner's gain over OLLA remains once the comparison includes rules calibrated on the learner's own cues?

Table~\ref{tab:lit} summarises eight closely related papers on two points that matter for that question. Across the \NLitN{} learning-based papers we screened, \NLitCalYes{} compared against a calibrated rule built on the learner's cues, \NLitCalPart{} did so partly, \NLitCalNo{} did not, and \NLitCalUnc{} could not be determined. These counts rest on abstracts and the text we could access, and an ``unclear'' entry is not a negative finding.

\begin{table}[t]
\caption{Eight closely related learned link-adaptation papers. ``Calibrated rule'': a non-learning rule calibrated on the learner's cues is a baseline. ``Parity'': the learner and that rule use comparable information. Entries come from the abstract and accessible text.}\label{tab:lit}
\centering\footnotesize
\setlength{\tabcolsep}{3pt}
\begin{tabular}{@{}p{1.55cm}p{2.5cm}p{1.8cm}p{0.9cm}p{0.8cm}@{}}
\toprule
Paper & Setting & HARQ model & Calibr.\ rule & Parity\\
\midrule
Saxena~\cite{saxena2022tuningfree} & single link, AWGN and fading & first transmission only & no & partial\\
Ye~\cite{ye2023drlla} & LTE/NR simulation & not stated & no & unclear\\
Kela~\cite{kela2022delayolla} & 5G UL, XR, system-level & system-level HARQ & no & no\\
Pang~\cite{pang2024td3olla} & LTE/NR simulation & not stated & no & unclear\\
You~\cite{you2026sim2real} & simulation and USRP/srsRAN testbed & parallel HARQ, ACK delay & partial & no\\
Wang (LOLLA)~\cite{wang2026lolla} & GPU NR stack, UL & immediate CRC, no HARQ delay & no & no\\
Demirel~\cite{demirel2025genrl} & system-level, multi-cell & system-level & partial & unclear\\
Wiesmayr (SALAD)~\cite{wiesmayr2025salad} & OAI, single user & real HARQ & partial & yes\\
\bottomrule
\end{tabular}
\end{table}

\subsection{Evaluation rigour}
Doherty et al.\ show that carefully tuned heuristics matched published RL results on optical-network problems~\cite{doherty2025hypehope}; this is the closest methodological precedent, in another domain. In recommender systems, Ferrari Dacrema et al.\ could reproduce only a minority of recent neural methods, and most of those were matched or beaten by simple, well-tuned heuristics~\cite{ferraridacrema2019progress}. Similar concerns about baselines, leakage and reproducibility are raised for learned systems and wireless machine learning in~\cite{mcgreivy2024weak,mangal2026gengap,kapoor2023leakage}. For RL specifically, variance across seeds and the reporting of intervals are discussed in~\cite{henderson2018matters,colas2018seeds}; Agarwal et al.\ recommend interval estimates, score distributions over runs and the probability of improvement instead of point estimates~\cite{agarwal2021precipice}, and we follow these recommendations. Pre-registration is advocated in~\cite{nosek2018prereg,albanie2022preregml}, and credibility of simulation studies in~\cite{pawlikowski2002credibility,kurkowski2005manet}. We apply these ideas to one concrete question in link adaptation.
